
\documentclass{article}
\usepackage{colortbl}
\usepackage{makecell}
\usepackage{multirow}
\usepackage{supertabular}

\begin{document}

\newcounter{utterance}

\centering \large Interaction Transcript for game `dond', experiment `coop\_en', episode 12 with Qwen3{-}30B{-}A3B{-}t0.0.
\vspace{24pt}

{ \footnotesize  \setcounter{utterance}{1}
\setlength{\tabcolsep}{0pt}
\begin{supertabular}{c@{$\;$}|p{.15\linewidth}@{}p{.15\linewidth}p{.15\linewidth}p{.15\linewidth}p{.15\linewidth}p{.15\linewidth}}
    \# & $\;$A & \multicolumn{4}{c}{Game Master} & $\;\:$B\\
    \hline

    \theutterance \stepcounter{utterance}  
    & & \multicolumn{4}{p{0.6\linewidth}}{
        \cellcolor[rgb]{0.9,0.9,0.9}{
            \makecell[{{p{\linewidth}}}]{
                \texttt{\tiny{[P1$\langle$GM]}}
                \texttt{You are playing a cooperative negotiation game in which you have to agree on how to divide a set of items among you and another player.} \\
\\ 
\texttt{Rules:} \\
\texttt{(a) You and the other player are given a set of items. Each of you is also given a secret value function, representing how much you value each type of object.} \\
\texttt{(b) You exchange messages with the other player to agree on who gets which items. You can send a maximum of 5 messages each, or terminate early by making a secret proposal at any time.} \\
\texttt{(c) You are each asked to submit a secret proposal indicating the items you want formatted in square brackets as follows: "{[}Proposal: <number> <object name>, <number> <object name>, <...>{]}"} \\
\texttt{(d) If your proposals are complementary, i.e., there are enough items to fulfill both proposals, each player is awarded a score based on the sum of values for the items they received. Otherwise, both of you get zero points.} \\
\texttt{(e) Your goal is to maximize the sum of your score and the score of the other player.} \\
\\ 
\texttt{Let us start.} \\
\\ 
\texttt{The set of available items is:} \\
\\ 
\texttt{3 rulers, 2 eggs, 1 battery} \\
\\ 
\texttt{Your secret value function is:} \\
\\ 
\texttt{ruler: 1, egg: 2, battery: 3} \\
\\ 
\texttt{IMPORTANT: Your messages, unless it is the secret proposal, are directly transmitted to the other player, so do not include any response to the rules or text announcing your message. To make a secret proposal, use the indicated format. Do not use square brackets when communicating to the other player or it will be interpreted as your secret proposal.} \\
            }
        }
    }
    & \\ \\

    \theutterance \stepcounter{utterance}  
    & \multicolumn{4}{p{0.6\linewidth}}{
        \cellcolor[rgb]{0.9,0.9,0.9}{
            \makecell[{{p{\linewidth}}}]{
                \texttt{\tiny{[P1$\rangle$GM]}}
                \texttt{<think>} \\
\texttt{Okay, let's see. I need to figure out how to divide the items between me and the other player to maximize the total score. The items are 3 rulers, 2 eggs, and 1 battery. My value function is ruler:1, egg:2, battery:3. But I don't know the other player's values. Hmm.} \\
\\ 
\texttt{First, since we're cooperating, the goal is to make sure that both of our proposals can be fulfilled. So I need to find a way to split the items so that neither of us is left out. But how?} \\
\\ 
\texttt{Let me think about the total value. If I take the battery, which is worth 3 to me, and maybe some eggs. The other player might value eggs more or less. Wait, but without knowing their values, I have to make an assumption. Maybe I should try to get the items that have higher value for me, but also leave enough for the other player.} \\
\\ 
\texttt{Alternatively, maybe I can propose a division that's fair. For example, if I take the battery and one egg, and the other player takes the rest. But I need to check if that's possible. Let's see: the battery is 1, so if I take that, and one egg (value 2), that's 3. The other player would get 3 rulers and 1 egg. But their value for rulers and eggs is unknown. But if their proposal can be fulfilled, then we both get points.} \\
\\ 
\texttt{Wait, but the problem says that if the proposals are complementary, meaning there are enough items to fulfill both. So I need to make sure that the items I want and the items they want don't overlap in a way that makes it impossible.} \\
\\ 
\texttt{But since I don't know their values, maybe I should try to propose a division that's likely to be compatible. For example, if I take the battery and one egg, and the other player takes the rulers and the remaining egg. But how do I communicate that?} \\
\\ 
\texttt{Alternatively, maybe I should send a message to the other player suggesting a split. But since I can only send 5 messages, I need to be efficient. Let me think: the other player might have different preferences. If I can get them to agree on a split where we each take items that are high value for us, but not conflicting.} \\
\\ 
\texttt{Wait, the secret proposal is the key. So maybe I should make a proposal that includes the battery and some} \\
            }
        }
    }
    & & \\ \\

    \theutterance \stepcounter{utterance}  
    & & & \multicolumn{2}{p{0.3\linewidth}}{
        \cellcolor[rgb]{0.9,0.9,0.9}{
            \makecell[{{p{\linewidth}}}]{
                \texttt{\tiny{[GM$|$GM]}}
                \texttt{continue} \\
            }
        }
    }
    & & \\ \\

    \theutterance \stepcounter{utterance}  
    & & & \multicolumn{2}{p{0.3\linewidth}}{
        \cellcolor[rgb]{0.9,0.9,0.9}{
            \makecell[{{p{\linewidth}}}]{
                \texttt{\tiny{[GM$|$GM]}}
                \texttt{<think>} \\
\texttt{Okay, let's see. I need to figure out how to divide the items between me and the other player to maximize the total score. The items are 3 rulers, 2 eggs, and 1 battery. My value function is ruler:1, egg:2, battery:3. But I don't know the other player's values. Hmm.} \\
\\ 
\texttt{First, since we're cooperating, the goal is to make sure that both of our proposals can be fulfilled. So I need to find a way to split the items so that neither of us is left out. But how?} \\
\\ 
\texttt{Let me think about the total value. If I take the battery, which is worth 3 to me, and maybe some eggs. The other player might value eggs more or less. Wait, but without knowing their values, I have to make an assumption. Maybe I should try to get the items that have higher value for me, but also leave enough for the other player.} \\
\\ 
\texttt{Alternatively, maybe I can propose a division that's fair. For example, if I take the battery and one egg, and the other player takes the rest. But I need to check if that's possible. Let's see: the battery is 1, so if I take that, and one egg (value 2), that's 3. The other player would get 3 rulers and 1 egg. But their value for rulers and eggs is unknown. But if their proposal can be fulfilled, then we both get points.} \\
\\ 
\texttt{Wait, but the problem says that if the proposals are complementary, meaning there are enough items to fulfill both. So I need to make sure that the items I want and the items they want don't overlap in a way that makes it impossible.} \\
\\ 
\texttt{But since I don't know their values, maybe I should try to propose a division that's likely to be compatible. For example, if I take the battery and one egg, and the other player takes the rulers and the remaining egg. But how do I communicate that?} \\
\\ 
\texttt{Alternatively, maybe I should send a message to the other player suggesting a split. But since I can only send 5 messages, I need to be efficient. Let me think: the other player might have different preferences. If I can get them to agree on a split where we each take items that are high value for us, but not conflicting.} \\
\\ 
\texttt{Wait, the secret proposal is the key. So maybe I should make a proposal that includes the battery and some} \\
            }
        }
    }
    & & \\ \\

    \theutterance \stepcounter{utterance}  
    & & \multicolumn{4}{p{0.6\linewidth}}{
        \cellcolor[rgb]{0.9,0.9,0.9}{
            \makecell[{{p{\linewidth}}}]{
                \texttt{\tiny{[GM$\rangle$P2]}}
                \texttt{You are playing a cooperative negotiation game in which you have to agree on how to divide a set of items among you and another player.} \\
\\ 
\texttt{Rules:} \\
\texttt{(a) You and the other player are given a set of items. Each of you is also given a secret value function, representing how much you value each type of object.} \\
\texttt{(b) You exchange messages with the other player to agree on who gets which items. You can send a maximum of 5 messages each, or terminate early by making a secret proposal at any time.} \\
\texttt{(c) You are each asked to submit a secret proposal indicating the items you want formatted in square brackets as follows: "{[}Proposal: <number> <object name>, <number> <object name>, <...>{]}"} \\
\texttt{(d) If your proposals are complementary, i.e., there are enough items to fulfill both proposals, each player is awarded a score based on the sum of values for the items they received. Otherwise, both of you get zero points.} \\
\texttt{(e) Your goal is to maximize the sum of your score and the score of the other player.} \\
\\ 
\texttt{Let us start.} \\
\\ 
\texttt{The set of available items is:} \\
\\ 
\texttt{3 rulers, 2 eggs, 1 battery} \\
\\ 
\texttt{Your secret value function is:} \\
\\ 
\texttt{ruler: 3, egg: 0, battery: 1} \\
\\ 
\texttt{IMPORTANT: Your messages, unless it is the secret proposal, are directly transmitted to the other player, so do not include any response to the rules or text announcing your message. To make a secret proposal, use the indicated format. Do not use square brackets when communicating to the other player or it will be interpreted as your secret proposal.} \\
\\ 
\\ 
\texttt{<think>} \\
\texttt{Okay, let's see. I need to figure out how to divide the items between me and the other player to maximize the total score. The items are 3 rulers, 2 eggs, and 1 battery. My value function is ruler:1, egg:2, battery:3. But I don't know the other player's values. Hmm.} \\
\\ 
\texttt{First, since we're cooperating, the goal is to make sure that both of our proposals can be fulfilled. So I need to find a way to split the items so that neither of us is left out. But how?} \\
\\ 
\texttt{Let me think about the total value. If I take the battery, which is worth 3 to me, and maybe some eggs. The other player might value eggs more or less. Wait, but without knowing their values, I have to make an assumption. Maybe I should try to get the items that have higher value for me, but also leave enough for the other player.} \\
\\ 
\texttt{Alternatively, maybe I can propose a division that's fair. For example, if I take the battery and one egg, and the other player takes the rest. But I need to check if that's possible. Let's see: the battery is 1, so if I take that, and one egg (value 2), that's 3. The other player would get 3 rulers and 1 egg. But their value for rulers and eggs is unknown. But if their proposal can be fulfilled, then we both get points.} \\
\\ 
\texttt{Wait, but the problem says that if the proposals are complementary, meaning there are enough items to fulfill both. So I need to make sure that the items I want and the items they want don't overlap in a way that makes it impossible.} \\
\\ 
\texttt{But since I don't know their values, maybe I should try to propose a division that's likely to be compatible. For example, if I take the battery and one egg, and the other player takes the rulers and the remaining egg. But how do I communicate that?} \\
\\ 
\texttt{Alternatively, maybe I should send a message to the other player suggesting a split. But since I can only send 5 messages, I need to be efficient. Let me think: the other player might have different preferences. If I can get them to agree on a split where we each take items that are high value for us, but not conflicting.} \\
\\ 
\texttt{Wait, the secret proposal is the key. So maybe I should make a proposal that includes the battery and some} \\
            }
        }
    }
    & \\ \\

    \theutterance \stepcounter{utterance}  
    & & & \multicolumn{4}{p{0.6\linewidth}}{
        \cellcolor[rgb]{0.9,0.9,0.9}{
            \makecell[{{p{\linewidth}}}]{
                \texttt{\tiny{[GM$\langle$P2]}}
                \texttt{<think>} \\
\texttt{Okay, let's try to figure out the best way to divide the items. The key is to maximize the total score, which is the sum of both players' scores. Since we don't know the other player's values, we have to make an assumption or find a split that's likely to work.} \\
\\ 
\texttt{First, let's list the items: 3 rulers, 2 eggs, 1 battery. My values are ruler: 3, egg: 0, battery: 1. Wait, no, the user said their secret value function is ruler: 3, egg: 0, battery: 1. Wait, no, the user's value function is ruler: 3, egg: 0, battery: 1? Wait, the problem says:} \\
\\ 
\texttt{"Your secret value function is:} \\
\\ 
\texttt{ruler: 3, egg: 0, battery: 1"} \\
\\ 
\texttt{Wait, that's confusing. So for me, rulers are worth 3, eggs 0, battery 1. So I value rulers the most, then battery, then eggs. But the other player's values are unknown. So I need to find a way to split the items so that both of us can get items that are valuable to us, and the proposals are complementary.} \\
\\ 
\texttt{But how? Since I don't know their values, maybe I should try to get the items that are high for me and leave the rest for them. But if they also value certain items, maybe we can find a split.} \\
\\ 
\texttt{For example, if I take the battery (value 1) and some rulers (value 3 each), and the other player takes the eggs (which I value 0, but maybe they value them more). But since I don't know their values, I have to guess.} \\
\\ 
\texttt{Alternatively, maybe the optimal split is to take the battery and one ruler, and the other player takes the rest. But I need to check if that's possible.} \\
\\ 
\texttt{Wait, the total items are 3 rulers, 2 eggs, 1 battery. If I take the battery and one ruler, that's 1 battery and 1 ruler. Then the other player could take 2 rulers and 2 eggs. But if their proposal is that, then it's complementary. But how do I know if that's what they want?} \\
\\ 
\texttt{Alternatively, maybe I should propose a split where I take the battery and one egg, but since I value eggs 0, that might not be} \\
            }
        }
    }
     \\ \\

    \theutterance \stepcounter{utterance}  
    & & & \multicolumn{2}{p{0.3\linewidth}}{
        \cellcolor[rgb]{0.9,0.9,0.9}{
            \makecell[{{p{\linewidth}}}]{
                \texttt{\tiny{[GM$|$GM]}}
                \texttt{continue} \\
            }
        }
    }
    & & \\ \\

    \theutterance \stepcounter{utterance}  
    & & & \multicolumn{2}{p{0.3\linewidth}}{
        \cellcolor[rgb]{0.9,0.9,0.9}{
            \makecell[{{p{\linewidth}}}]{
                \texttt{\tiny{[GM$|$GM]}}
                \texttt{<think>} \\
\texttt{Okay, let's try to figure out the best way to divide the items. The key is to maximize the total score, which is the sum of both players' scores. Since we don't know the other player's values, we have to make an assumption or find a split that's likely to work.} \\
\\ 
\texttt{First, let's list the items: 3 rulers, 2 eggs, 1 battery. My values are ruler: 3, egg: 0, battery: 1. Wait, no, the user said their secret value function is ruler: 3, egg: 0, battery: 1. Wait, no, the user's value function is ruler: 3, egg: 0, battery: 1? Wait, the problem says:} \\
\\ 
\texttt{"Your secret value function is:} \\
\\ 
\texttt{ruler: 3, egg: 0, battery: 1"} \\
\\ 
\texttt{Wait, that's confusing. So for me, rulers are worth 3, eggs 0, battery 1. So I value rulers the most, then battery, then eggs. But the other player's values are unknown. So I need to find a way to split the items so that both of us can get items that are valuable to us, and the proposals are complementary.} \\
\\ 
\texttt{But how? Since I don't know their values, maybe I should try to get the items that are high for me and leave the rest for them. But if they also value certain items, maybe we can find a split.} \\
\\ 
\texttt{For example, if I take the battery (value 1) and some rulers (value 3 each), and the other player takes the eggs (which I value 0, but maybe they value them more). But since I don't know their values, I have to guess.} \\
\\ 
\texttt{Alternatively, maybe the optimal split is to take the battery and one ruler, and the other player takes the rest. But I need to check if that's possible.} \\
\\ 
\texttt{Wait, the total items are 3 rulers, 2 eggs, 1 battery. If I take the battery and one ruler, that's 1 battery and 1 ruler. Then the other player could take 2 rulers and 2 eggs. But if their proposal is that, then it's complementary. But how do I know if that's what they want?} \\
\\ 
\texttt{Alternatively, maybe I should propose a split where I take the battery and one egg, but since I value eggs 0, that might not be} \\
            }
        }
    }
    & & \\ \\

    \theutterance \stepcounter{utterance}  
    & & \multicolumn{4}{p{0.6\linewidth}}{
        \cellcolor[rgb]{0.9,0.9,0.9}{
            \makecell[{{p{\linewidth}}}]{
                \texttt{\tiny{[P1$\langle$GM]}}
                \texttt{<think>} \\
\texttt{Okay, let's try to figure out the best way to divide the items. The key is to maximize the total score, which is the sum of both players' scores. Since we don't know the other player's values, we have to make an assumption or find a split that's likely to work.} \\
\\ 
\texttt{First, let's list the items: 3 rulers, 2 eggs, 1 battery. My values are ruler: 3, egg: 0, battery: 1. Wait, no, the user said their secret value function is ruler: 3, egg: 0, battery: 1. Wait, no, the user's value function is ruler: 3, egg: 0, battery: 1? Wait, the problem says:} \\
\\ 
\texttt{"Your secret value function is:} \\
\\ 
\texttt{ruler: 3, egg: 0, battery: 1"} \\
\\ 
\texttt{Wait, that's confusing. So for me, rulers are worth 3, eggs 0, battery 1. So I value rulers the most, then battery, then eggs. But the other player's values are unknown. So I need to find a way to split the items so that both of us can get items that are valuable to us, and the proposals are complementary.} \\
\\ 
\texttt{But how? Since I don't know their values, maybe I should try to get the items that are high for me and leave the rest for them. But if they also value certain items, maybe we can find a split.} \\
\\ 
\texttt{For example, if I take the battery (value 1) and some rulers (value 3 each), and the other player takes the eggs (which I value 0, but maybe they value them more). But since I don't know their values, I have to guess.} \\
\\ 
\texttt{Alternatively, maybe the optimal split is to take the battery and one ruler, and the other player takes the rest. But I need to check if that's possible.} \\
\\ 
\texttt{Wait, the total items are 3 rulers, 2 eggs, 1 battery. If I take the battery and one ruler, that's 1 battery and 1 ruler. Then the other player could take 2 rulers and 2 eggs. But if their proposal is that, then it's complementary. But how do I know if that's what they want?} \\
\\ 
\texttt{Alternatively, maybe I should propose a split where I take the battery and one egg, but since I value eggs 0, that might not be} \\
            }
        }
    }
    & \\ \\

    \theutterance \stepcounter{utterance}  
    & \multicolumn{4}{p{0.6\linewidth}}{
        \cellcolor[rgb]{0.9,0.9,0.9}{
            \makecell[{{p{\linewidth}}}]{
                \texttt{\tiny{[P1$\rangle$GM]}}
                \texttt{<think>} \\
\texttt{Okay, let's try to approach this step by step. The goal is to maximize the total score, which is the sum of both players' scores. Since I don't know the other player's value function, I need to make an educated guess based on the items and my own values.} \\
\\ 
\texttt{My value function is: rulers (3), eggs (0), battery (1). So I value rulers the most, then battery, then eggs. The other player's values are unknown, but they might have different priorities. To maximize the total, I should aim for a division where both of us get items that are valuable to us, and the proposals don't conflict.} \\
\\ 
\texttt{Let's consider possible splits:} \\
\\ 
\texttt{1. **I take the battery (value 1) and some rulers.** Since rulers are high value for me, taking a couple might be good. But if I take too many, the other player might not get enough. However, since I don't know their values, maybe I should take the battery and one ruler, leaving 2 rulers, 2 eggs, and 0 battery for the other player. But if they value eggs or rulers more, this could work.} \\
\\ 
\texttt{2. **I take the battery and eggs.** But eggs are worth 0 to me, so that's not efficient. Unless the other player values eggs, but I can't assume that.} \\
\\ 
\texttt{3. **I take all rulers.** That would give me 3*3=9, but the other player would get nothing, leading to zero points. Not good.} \\
\\ 
\texttt{Alternatively, maybe a fair split where both get some high{-}value items. Since I value rulers the most, maybe I take 2 rulers (value 6), and the other player takes 1 ruler, 2 eggs, and the battery. But again, without knowing their values, this is speculative.} \\
\\ 
\texttt{Wait, the key is that the proposals must be complementary. So if I propose a certain set, the other player's proposal must not require more items than available. So I need to make a proposal that leaves enough items for them.} \\
\\ 
\texttt{Let me think of possible proposals:} \\
\\ 
\texttt{{-} If I propose {[}battery, ruler{]}, that's 1 battery and 1 ruler. Then the other player could propose {[}ruler, ruler, egg, egg{]}, which would use the remaining 2 rulers and 2 eggs. That's complementary. But if they value rulers and eggs, this could work. But since I don't} \\
            }
        }
    }
    & & \\ \\

    \theutterance \stepcounter{utterance}  
    & & & \multicolumn{2}{p{0.3\linewidth}}{
        \cellcolor[rgb]{0.9,0.9,0.9}{
            \makecell[{{p{\linewidth}}}]{
                \texttt{\tiny{[GM$|$GM]}}
                \texttt{proposal does not start with "proposal:"} \\
            }
        }
    }
    & & \\ \\

\end{supertabular}
}

\end{document}
